\documentclass[aspectratio=169,10pt,english]{beamer}

\input{../../../_preambulo_beamer.tex}
\input{../../../_comandos_beamer.tex}

\selectlanguage{english}

\title{MA1001B - Statistical Modeling for Decision Making}
\subtitle{Lesson 26: Point Estimation: Unbiasedness, Efficiency, Consistency}
\author{}
\institute{Tecnol\'ogico de Monterrey}
\date{}

\begin{document}

\begin{frame}[plain]
  \vspace{1.2cm}
  {\Large\bfseries MA1001B - Statistical Modeling for Decision Making\par}
  \vspace{0.7cm}
  {\large Lesson 26: Point Estimation Properties\par}
  \vspace{0.7cm}
  {\color{TecRojo}\rule{\textwidth}{1.2pt}}
  \vspace{0.8cm}

  MA1001B Faculty\\[0.3cm]
  Term\\[0.3cm]
  Tecnol\'ogico de Monterrey
\end{frame}

\begin{frame}{The Estimator Question}
  \begin{alertblock}{Guiding question}
    If costs are $N(100,15)$, why prefer $\bar x$ over a single observation
    $X_1$ as a point estimator of $\mu$?
  \end{alertblock}

  \vspace{0.6em}
  By the end of this lesson, you can:
  \begin{itemize}
    \item define unbiasedness as $E[\hat\theta]=\theta$;
    \item compare efficiency through estimator variances;
    \item recognize consistency as spread shrinking with $n$;
    \item explain why unbiasedness does not imply a small sample is precise.
  \end{itemize}

  \vfill
  \scriptsize All costs used today are synthetic. No real company data are used.
\end{frame}

\begin{frame}{Two Unbiased Estimators}
  \small
  Population $X\sim N(100,15^2)$. Both
  \[
    E[\bar x]=\mu\qquad\text{and}\qquad E[X_1]=\mu.
  \]
  8{,}000 seed-42 samples of $n=20$:
  \[
    \text{mean of }\bar x=99.957310,\quad \text{bias}=-0.042690,
  \]
  \[
    \text{mean of }X_1=99.916074,\quad \text{bias}=-0.083926.
  \]

  \vspace{0.3em}
  \begin{exampleblock}{Unbiasedness}
    Correct on average is not the same as small variance.
  \end{exampleblock}
\end{frame}

\begin{frame}[fragile]{Step 1: Unbiasedness}
  \begin{columns}[T]
    \begin{column}{0.42\textwidth}
      \small
      \begin{lstlisting}
xbar = draws.mean(axis=1)
one = draws[:, 0]
bias_xbar = xbar.mean() - 100
      \end{lstlisting}

      \begin{block}{Verified output}
        Mean of $\bar x=99.957310$\\
        Mean of $X_1=99.916074$
      \end{block}
    \end{column}
    \begin{column}{0.58\textwidth}
      \centering
      \includegraphics[width=\textwidth]{../figures/point_estimation_01_unbiasedness.png}
      \vspace{0.2em}
      \scriptsize Sampling means from Step 1
    \end{column}
  \end{columns}
\end{frame}

\begin{frame}{Efficiency and Consistency}
  \small
  At $n=20$,
  \[
    \mathrm{Var}(X_1)=225,\qquad \mathrm{Var}(\bar x)=11.25.
  \]
  Empirical values: $227.612313$ versus $11.242349$, ratio $20.245975$.
  Raising $n$ to 200 shrinks $\mathrm{SD}(\bar x)$ to $1.052019$ versus
  formula $15/\sqrt{200}=1.060660$.
\end{frame}

\begin{frame}[fragile]{Step 2: Spread Comparison}
  \begin{columns}[T]
    \begin{column}{0.40\textwidth}
      \small
      \begin{lstlisting}
var_xbar = xbar.var()
var_one = one.var()
ratio = var_one / var_xbar
      \end{lstlisting}

      \begin{block}{Verified output}
        $\mathrm{SD}(\bar x_{20})=3.352961$\\
        $\mathrm{SD}(\bar x_{200})=1.052019$
      \end{block}
    \end{column}
    \begin{column}{0.60\textwidth}
      \centering
      \includegraphics[width=\textwidth]{../figures/point_estimation_02_efficiency_consistency.png}
      \vspace{0.2em}
      \scriptsize Efficiency and consistency from Step 2
    \end{column}
  \end{columns}
\end{frame}

\begin{frame}{Step 3: Unbiased but Too Noisy}
  \begin{columns}[T]
    \begin{column}{0.42\textwidth}
      \small
      At $n=5$, $\bar x$ is still unbiased:
      mean $100.066209$, SD $6.737762$.

      \vspace{0.4em}
      $P(|\bar x-100|>10)=0.137875$.\\
      At $n=200$ that event is $0.000000$.

      \vspace{0.4em}
      \textbf{Limit:} unbiasedness does not imply small variance.
      Lesson 27 wraps $\bar x$ in a $z$ interval.
    \end{column}
    \begin{column}{0.58\textwidth}
      \centering
      \includegraphics[width=\textwidth]{../figures/point_estimation_03_small_sample_variance.png}
      \vspace{0.2em}
      \scriptsize Small-sample noise from Step 3
    \end{column}
  \end{columns}
\end{frame}

\begin{frame}{Concept Check}
  \begin{enumerate}
    \item Why are both $\bar x$ and $X_1$ unbiased for $\mu$?
    \item What does the ratio $\mathrm{Var}(X_1)/\mathrm{Var}(\bar x)\approx 20$
      say at $n=20$?
    \item How does $\mathrm{SD}(\bar x)$ change from $n=20$ to $n=200$?
    \item Why can an unbiased $n=5$ mean still be the wrong decision tool?
  \end{enumerate}

  \vspace{0.6em}
  \begin{block}{Connection to Lesson 27}
    The session can continue directly into a $z$ confidence interval for a
    mean when $\sigma$ is known.
  \end{block}

  \vfill
  \scriptsize Source: OpenStax, \textit{Introductory Business Statistics 2e},
  Chapter 8 introduction. CC BY-NC-SA.
\end{frame}

\appendix



\end{document}
